\section{Introduction}

\textit{Data-Driven computational mechanics} (DDCM) represents the latest paradigm in continuum mechanics, which can, of course, be employed to a large range of applications. For instance, in the context of the diffusion-reaction problem, the gradient of a concentration scalar field, corresponding to a conserved quantity, and the flux vector field over a physical domain are then to be found as the closest to some field extension of a given finite set of empirically determined gradient-flux pairs. Such a problem can be stated as an optimization problem in which discrete and continuous fields are to be simultaneously sought such that the physical conservation law and the compability condition are satisfied in some weak sense. In this context, a more basic question on the existence of the continuous fields provided that the discrete fields are fixed and given \textit{a priori} naturally arises. This ``optimization subproblem'' is actually a continuous optimization problem that can, in principle, be treated computationally and analyzed using standard tools from the classical theory of linear \textit{partial differential equations} (PDEs). Therefore, answering that question represents the main target of the present contribution. To place the proposed ideas into an appropriate setting, we provide below some of the most important contributions to present date, highlight the difference of the present work with respect to existing literature and state the novelty from the methodological point of view.

DDCM is a very active field of research whose first contributions can be traced back to less than a decade ago and has mostly been developed within solid/structural mechanics. For instance, this principle is applied to the static analysis of multi-bar structures \cite{Kirchdoerfer2016}. Based on similar ideas, a data-driven magnetostatic solver is proposed in \cite{Galetzka2020}. In the following, this initial approach is modified to improve robustness due to the presence of outliers in the constitutive data \cite{Kirchdoerfer2017}. Finally, this is extended to the dynamic analysis of structures that exhibit configuration-independent mass matrices \cite{Kirchdoerfer2018}. For the sake of completeness, we must mention that there are some efforts addressing the existence of solutions for data-driven elasticity \cite{Conti2018,Conti2020}. There is also a hybrid approach that relies on a smooth constitutive manifold reconstructed from the data at a off-line pre-step. This principle is applied to time-independent problems \cite{Gebhardt:2020a}. Then, the approach is also modified to handle time-dependent problems \cite{Gebhardt:2020b}. Lastly, this is extended to deal with \textit{mathematical programs with complementarity constraints} (MPCC) such as those arising in contact mechanics  \cite{Gebhardt2024a}. Recently in  \cite{Gebhardt2024b}, the structural properties of the DDCM problem discretized by means of the \textit{finite element method} (FEM) is investigated and its global solvability is established. Remarkably, there are neither antecedents of DDCM in the context of the diffusion-reaction problem nor direct connections to PDE-regularized methods to analyze empirical data \cite{Ponti2021} or to inhomogeneities identification from boundary measurements \cite{Ammari2004}.

All the works introduced above (that deal with the numerical solution of the discrete-continuous optimization problem in DDCM) require along the solving process to compute the solution of a type of continuous optimization subproblem that is a particular case of the optimization problem addressed in the present work. This is due to that those works adopt an \textit{alternating direction method} (ADM). This work can be considered then as a first step towards the application of DDCM in the context of the diffusion-reaction problem, which to the best of our knowledge has not been realized so far. In previous works, where once the assignation of the data pairs has been done, the optimization subproblem provides two sets of optimality conditions that are totally decoupled and can be solved independently. In the present work, we consider a reaction term present in the physical conservation law that automatically couples all the optimality conditions and thus, the well-posedness of this rather more general optimization problem ought to be established right from the beginning. Moreover, our analysis considers two- and three spatial dimensions in contrast to those more practical contributions introduced before that focus on problems defined on a closed interval of the real line \cite{Gebhardt:2020a, Gebhardt:2020b, Gebhardt2024a, Gebhardt2024b}. In addition, we also introduce a new finite element formulation, establish its well-posedness for the chosen finite element spaces and show its consistency with respect to the continuous formulations posed in the setting of infinite-dimensional Hilbert spaces. The desired stability properties are achieved by directly approximating the saddle-point problem, approach that exhibits some connection to other well-established approaches \cite{Masud2002, Bochev2006, Badia2010, Burman2023}. Moreover, we also consider the case in which the data sets are generated from a Darcy-Forcheimer flow type \cite{Girault2008}.

The remaining of this article is organized as follows. In Section 2, we describe the optimization problem in the reaction-diffusion context. Relying on the theory for quadratic optimization problems subject to linear constraints, we prove that this optimization problem in its continuous setting is indeed solvable with a unique solution. Moreover, we also introduce its spatial approximation by adopting a finite element approach. Then, we prove that the proposed approximating problem is also well posed and consistent. In Section 3, we introduce the numerical implementation of the spatially discretized problem, which is realized by its coding within the Firedrake finite element platform. The resulting code is then briefly explained and made publicly available. This development is followed as well by a series of examples, of increasing difficulty, that illustrate the remarkable properties of the setting proposed. Finally, in Section 4, we conclude with remarks, potential limitations and possible future work.  
